\documentclass[10pt]{article}
\usepackage[margin=0.9in]{geometry}
\usepackage{amsmath,amsthm,booktabs}
\usepackage[hidelinks]{hyperref}
\DeclareMathSizes{10}{10}{7}{7}
\newtheorem{theorem}{Theorem}
\newtheorem{lemma}[theorem]{Lemma}
\newcommand{\gp}{\mathbin{\vcenter{\hbox{\setlength{\fboxsep}{0pt}%
  \fbox{\rule{0pt}{.51em}\rule{.51em}{0pt}}}}}}
\title{Vizing's domination inequality for pairs\\of graphs of order at most eleven}
\author{Claude Make No Mistakes}
\date{26 November 2026}
\begin{document}
\maketitle

\begin{abstract}
We verify Vizing's inequality for every pair of finite simple graphs of order
at most eleven. An exhaustive enumeration reduces the test to 63 possible
factors. Rational weights on subsets of each factor certify 1,971 of the
2,016 unordered pairs of these factors, including diagonal pairs. We verify
the remaining 45 pairs by generating and checking DRAT refutations. The
same weights give a lower bound of \(61/66\) times the conjectured product
when one factor has order at most eleven and the other has arbitrary order.
\end{abstract}

Write \(N_G[S]\) for the union of the closed neighborhoods of \(S\) in a
finite simple graph \(G\). The domination number \(\gamma(G)\) is the minimum
size of a set \(S\) with \(N_G[S]=V(G)\). Vizing conjectured that
\[
    \gamma(G\gp H)\ \ge\ \gamma(G)\gamma(H)
\]
for every pair of such graphs. In the Cartesian product, adjacent vertices
agree in one coordinate and are adjacent in the other. The September 2026
revision by Aliabadi and Krop gives a universal factor of \(2/3\) \cite{ak}.

\begin{theorem}
Let \(G,H\) be finite simple graphs.
\begin{enumerate}
\item If both orders are at most eleven, then
      \(\gamma(G\gp H)\ge\gamma(G)\gamma(H)\).
\item If \(|V(G)|\le 9\), the same inequality holds for every \(H\).
      It also holds if \(|V(G)|\le 10\) and \(\gamma(H)\le 4\).
\item If \(|V(G)|\le 11\), then
      \(\gamma(G\gp H)\ge(61/66)\gamma(G)\gamma(H)\).
      If \(|V(G)|\le 10\), the factor improves to \(19/20\).
\end{enumerate}
\end{theorem}

\section{Reducing the factors}

An empty factor makes the inequality immediate. Domination numbers add over
components, and the product distributes over components of either factor.
Thus a violation has a connected violating pair within the same order bounds.
For a fixed \(H\) and a fixed target coefficient \(0<c\le1\), choose a connected
factor \(G\) of smallest order with
\(\gamma(G\gp H)<c\gamma(G)\gamma(H)\). Among these, maximize the edge count.
We apply two reductions from Bre\v{s}ar et al.\ \cite[\S 6]{survey}.

Adding a missing edge of \(G\) can only decrease the product domination number.
Hence the addition must strictly decrease \(\gamma(G)\). Identifying any
two distinct vertices of \(G\) and deleting any resulting loops also gives a
product with domination number at most \(\gamma(G\gp H)\). The quotient map
times the identity on \(H\) is surjective and takes every closed neighborhood
into the closed neighborhood of its image. It therefore sends a dominating
set of the original product to one in the quotient product. The identification
must strictly decrease \(\gamma(G)\), since it reduces the order and preserves
connectedness. For a counterexample within a bound on both orders, choose
a pair minimizing the sum of orders and then maximizing the sum of edge
counts. The same argument applies to each factor in that pair.

Let \(\rho(G)\) be the maximum size of a collection of vertices with pairwise
disjoint closed neighborhoods. It is at most \(\gamma(G)\). The known exact
inequality for arbitrary partners applies when \(\gamma(G)\le3\), or when
\(\gamma(G)\le\rho(G)+1\) \cite{lowgamma,half,survey}. The case
\(\gamma=\rho\) also follows by projecting a product dominating set from
each of the disjoint neighborhoods onto \(H\). Consequently every selected
factor must pass all four tests in Table~\ref{tab:scan}. These tests apply
to each target coefficient in the theorem.

\begin{table}[h]
\centering
\begin{tabular}{@{}rrrrrr@{}}
\toprule
Order & Connected & \(\gamma\ge4\) & \(\rho\le\gamma-2\) & Edge additions & Identifications\\
\midrule
1--7 & 996 & 0 & 0 & 0 & 0\\
8 & 11,117 & 6 & 0 & 0 & 0\\
9 & 261,080 & 191 & 0 & 0 & 0\\
10 & 11,716,571 & 9,665 & 74 & 5 & 5\\
11 & 1,006,700,565 & 738,353 & 7,068 & 76 & 58\\
\bottomrule
\end{tabular}
\caption{Successive filters on the connected graphs generated by nauty 2.9.3 \cite{nauty}.
The last two columns require that every missing-edge addition and every
pair identification, respectively, decrease the original domination number.}
\label{tab:scan}
\end{table}

The generator \texttt{geng -cq} supplies the connected unlabeled graphs.
An output hook tests closed-neighborhood unions, disjointness and the
modified graphs using integer bit sets. A second run prunes an induced
prefix on \(n-r\) vertices if it has a dominating set of size at most \(3-r\),
since adding the \(r\) future vertices would give a dominating set of size
at most three in every completion. A third run at order eleven additionally
uses maximum degree at most six. In general, if a vertex leaves at most two
vertices outside its closed neighborhood, \(\gamma\le3\). If the vertex
leaves exactly three, then \(\gamma\le4\), and equality forces the three
remaining vertices to have pairwise disjoint closed neighborhoods. Otherwise
one vertex covers two of the three and yields a dominating triple.
Thus every graph passing the first two filters has maximum degree at most
\(n-5\). The runs agree on the complete sets of retained graph6 records.
An independent streaming implementation agrees with the unpruned counts
through order ten and rechecks every retained order-eleven factor.
All 63 retained factors have \((\gamma,\rho)=(4,2)\).

\section{Weighted subset certificates}

The following bound extends the counting condition for a fair reception
\cite{reception}. Blocks may intersect. For a nonempty block \(B\subseteq V(G)\)
and \(S\subseteq V(G)\), say \(B\) is active at \(S\) if
\(B\cap S\ne\emptyset\) or \(B\subseteq N_G[S]\).
Let \(\mathcal A(S)\) be the collection of active blocks.

\begin{lemma}
Give the nonempty blocks \(B\) nonnegative integer weights \(w_B\).
Suppose that for an integer \(q>0\) and every \(S\subseteq V(G)\),
\[
    \sum_{B\in\mathcal A(S)}w_B \ \le\ q|S|.
\]
Let \(W=\sum_B w_B\). For every \(H\) one has
\(\gamma(G\gp H)\ge (W/q)\gamma(H)\).
\end{lemma}

\begin{proof}
Let \(D\) dominate \(G\gp H\) and \(S_h=\{g:(g,h)\in D\}\).
Put \(U_B=\{h:B\text{ is active at }S_h\}\).
If \(h\notin U_B\), choose \(g\in B\setminus N_G[S_h]\).
The vertex \((g,h)\) must be dominated at \((g,h')\) for some neighbor
\(h'\) of \(h\). Then \(g\in B\cap S_{h'}\), so \(h'\in U_B\).
Thus \(U_B\) dominates \(H\), and summing the assumed capacity over \(h\)
gives \(q|D|\ge\sum_B w_B|U_B|\ge W\gamma(H)\).
\end{proof}

The ledger stores rational weights found by linear programming and checked
for every relevant vertex subset by integer arithmetic. Subsets with
\(q|S|\ge W\) satisfy the condition automatically. The minimum checked
weight is \(19/5\) over the order-ten factors and \(122/33\) over the order-eleven
factors. Together with the minimal-factor argument, these weights prove
part 3. They also prove the exact order-ten assertion for \(\gamma(H)\le4\)
by rounding up, since \((19/5)\gamma(H)>4\gamma(H)-1\) when \(1\le\gamma(H)\le4\).
There are no factors to check through order nine. Eleven-vertex factors
with weight \(>15/4\) satisfy the exact inequality whenever their partner
has domination number four. Nine factors have checked weight
at most \(15/4\), leaving 45 unordered pairs to check.

\section{Checking the remaining pairs}

For a fixed pair \(G,H\), set \(k=\gamma(G)\), \(\ell=\gamma(H)\) and
\(b=k\ell-1\). Boolean variables select vertices of the Cartesian product.
For each product vertex, a clause requires some selected vertex in its
closed neighborhood. A sequential counter requires at most \(b\) selected
vertices. The weighted certificates give additional necessary conditions.
If \(D\) satisfies the budget, let \(s=qb-\ell W\). Then
\[
\begin{split}
q|D|-\ell W
={}&\sum_{h\in V(H)}\left(q|S_h|-\sum_{B\in\mathcal A(S_h)}w_B\right)\\
 &+\sum_B w_B\bigl(|U_B|-\ell\bigr) \ \le\ s.
\end{split}
\]
Each summand is nonnegative. A fiber with vertex set \(S\) can occur only
if its deficiency \(q|S|-\sum_{B\in\mathcal A(S)}w_B\) is at most \(s\);
in particular \(|S|\le\lfloor(W+s)/q\rfloor\).
For \(w_B>0\) we also have
\(|U_B|\le\ell+\lfloor s/w_B\rfloor\).

The formula excludes each over-budget fiber size and each smaller vertex set
whose deficiency exceeds \(s\). Activation of a block at a permitted fiber
forces an auxiliary flag. A separate sequential counter bounds the number
of flags for that block. These clauses are added in both product directions.
Any dominating set of size at most \(b\) extends to a satisfying assignment
by taking the flags to be exactly the activations and extending the sequential
counters. In some cases we also use that the intersection of an empty row
and an empty column would be undominated. We impose lexicographic order
against checked product automorphisms. Every orbit of dominating sets
contains a lexicographically least member satisfying these constraints.
Conversely the base clauses impose product domination and the budget.
CaDiCaL 1.9.5 emits a DRAT refutation for each of the 45 formulas.
DRAT-trim checks each trace \cite{drat}. They can be regenerated and replayed from the
recorded graph6 pairs and weights. This completes part 1.

\section{Literature and search scope}

An August 2026 preprint claims a Lean resolution \cite{reed}. The final
theorem in \texttt{VizingDescentProof.lean} takes Vizing's inequality
for the pair as a hypothesis. The preprint also asserts
multiplicativity of ordinary fractional domination under Cartesian products.
For \(K_2\gp K_2=C_4\) the factor fractional domination numbers are both
one and the product value is \(4/3\). Indeed the four closed-neighborhood
constraints count each vertex three times, and the constant \(1/3\)
meets all of them with equality. The formal conclusion is conditional,
and this four-vertex example falsifies the fractional identity.

A separate search generated 22 critical factors of orders 13 through 16
with domination number five. The bounded SAT campaign tested their unordered
pairs together with the order-ten factor of weight \(19/5\). It found no
counterexample. Of the 276 runs, five returned UNSAT and 271 reached their
conflict limits. These runs lie outside the theorem and their solver
verdicts were recorded without DRAT replay.

\begin{thebibliography}{9}
\bibitem{ak}
M.~Aliabadi and E.~Krop,
\emph{A \(2/3\) Bound for Vizing's Conjecture},
\href{https://arxiv.org/abs/2607.01109v3}{arXiv:2607.01109v3}, 2026.
\bibitem{survey}
B.~Bre\v{s}ar et al.,
\emph{Vizing's conjecture: a survey and recent results},
\href{https://doi.org/10.1002/jgt.20565}{J. Graph Theory 69 (2012), 46--76}.
\bibitem{lowgamma}
B.~Bre\v{s}ar, \emph{Vizing's Conjecture for Graphs with Domination Number 3:
a New Proof},
\href{https://doi.org/10.37236/5182}{Electron. J. Combin. 22(3) (2015), \#P3.38}.
\bibitem{half}
B.~L.~Hartnell and D.~F.~Rall,
\emph{Vizing's conjecture and the one-half argument},
\href{https://doi.org/10.7151/dmgt.1018}{Discuss. Math. Graph Theory 15 (1995), 205--216}.
\bibitem{reception}
B.~Bre\v{s}ar and D.~F.~Rall,
\emph{Fair reception and Vizing's conjecture},
\href{https://doi.org/10.1002/jgt.20366}{J. Graph Theory 61 (2009), 45--54}.
\bibitem{nauty}
B.~D.~McKay and A.~Piperno, \emph{Practical graph isomorphism, II},
\href{https://doi.org/10.1016/j.jsc.2013.09.003}{J. Symbolic Comput. 60 (2014), 94--112}.
\bibitem{drat}
N.~Wetzler, M.~J.~H.~Heule, and W.~A.~Hunt Jr.,
\emph{DRAT-trim: Efficient Checking and Trimming Using Expressive Clausal Proofs},
\href{https://doi.org/10.1007/978-3-319-09284-3_31}{SAT 2014, 422--429}.
\bibitem{reed}
J.~Reed, \emph{Vizing's Impossible Descent},
\href{https://zenodo.org/records/21854681}{Zenodo record 21854681}, 2026.
\end{thebibliography}
\end{document}
